\begin{textsample}{POS Dim 1 – 2011\_c4 – Score -3.00 – 2011\_c4/enade\_2011\_tec\_constr\_f1a5d330d7.txt}  \label{ex:f1_pos_325}
CAPÍTULO 4 PERCEPÇÃO DA PROVA As análises feitas neste capítulo tratam das percepções dos concluintes da Área de Tecnologia em Construção de Edifícios sobre a prova aplicada no ENADE/2011.
Estas percepções foram mensuradas por meio de nove questões que avaliaram desde o grau de dificuldade da prova até o tempo gasto para concluí -la.
As percepções sobre a prova foram relacionadas com o desempenho dos estudantes e com a Grande Região de funcionamento do curso.
O questionário de percepção da prova encontra -se ao final do Anexo IV, que traz a reprodução da prova.
O desempenho dos estudantes foi classificado em quatro quartos.
Para tanto, esse desempenho foi ordenado de forma ascendente.
O percentil 25, P25, também conhecido como primeiro quartil, é a nota de desempenho que deixa um quarto ( 25 \% ) dos valores observados abaixo e três quartos acima.
A Figura 1 apresenta uma ilustração deste conceito.
O quarto inferior de desempenho é composto pelas notas abaixo do primeiro quartil.
Já o percentil 75, P75, também conhecido como terceiro quartil, é o valor para o qual há três quartos ( 75 \% ) dos dados abaixo e um quarto acima dele.
O quarto superior de desempenho é composto pelas notas iguais ou acima do terceiro quartil.
O percentil 50, P50, também conhecido como mediana, é o valor que divide as notas em dois conjuntos de igual tamanho.
O segundo quarto inclui valores entre o primeiro quartil ( P25 ) e a mediana.
O terceiro quarto contém os valores entre a mediana ( P50 ) e o terceiro quartil ( P75 ).
Vale ressaltar que percentis, quartis e medianas são pontos que não obrigatoriamente pertencem ao conjunto original de dados, ao passo que os quartos são subconjuntos dos dados originais.
A seguir, serão apresentados gráficos com resultados selecionados, relativos às nove questões avaliadas por grupos de estudantes.
Os gráficos apresentam nas barras o percentual de alunos que assinalaram uma das opções ou a soma das porcentagens daqueles que assinalaram duas ( ou três ) delas.
Por exemplo, para as questões 1 e 2, os gráficos apresentam a porcentagem total de participantes que assinalaram as opções ( D ) difícil e ( E ) muito difícil.
Em cada barra foram assinalados também os extremos do intervalo de confiança de 95 \% como linhas verticais unidas por uma linha horizontal na forma da letra H maiúscula, semelhantemente aos gráficos do Capítulo 3.
As Tabelas no Anexo II apresentam os valores absolutos e a distribuição percentual das alternativas válidas das nove questões, segundo o mesmo recorte de desempenho dos alunos e Grande Região de funcionamento do curso. 4.1 GRAU DE DIFICULDADE DA PROVA 4.1.1 Componente de Formação Geral Ao avaliarem “ Qual o grau de dificuldade desta prova na parte de Formação Geral? ” ( Questão 1 ), 15,0 \% do grupo de inscritos e presentes optaram pelas alternativas difícil ou muito difícil.
Entretanto, para mais da metade dos estudantes ( 68,9 \% ), o Componente de Formação Geral da prova foi considerado com grau de dificuldade médio ( Gráfico 4.1, Gráfico 4.2 e, no Anexo II, a Tabela II.1 ).
O percentual de estudantes que consideraram a prova como difícil ou muito difícil foi maior na região Norte, onde a proporção foi de 22,2 \%, enquanto a de menor incidência foi a Centro-Oeste, com 8,7 \%.
No Gráfico 4.1 é possível observar que a diferença entre as regiões não é estatisticamente significativa.
Nas Grandes Regiões, a proporção de presentes à prova que consideraram o Componente de Formação Geral como sendo de grau de dificuldade médio esteve entre 57,1 \% na região Sul e 82,6 \% na região Centro-Oeste.
O percentual de alunos que consideraram a prova difícil ou muito difícil foi semelhante no 1º e no 3º quarto de desempenho.
Verifica -se uma tendência decrescente conforme o aumento do desempenho dos estudantes a partir do 2º quarto.
Não existem diferenças estatisticamente significativas entre os quartos de desempenho.
A alternativa modal para esta pergunta foi médio. 4.1.2 Componente de Conhecimento Específico Ao responderem à Questão 2 — “ Qual o grau de dificuldade desta prova na parte de Componente Específico? ” — 21,2 \% do grupo de estudantes classificaram -na como difícil ou muito difícil.
Além disso, o Componente de Conhecimento Específico da prova foi considerado com grau de dificuldade médio por 67,5 \% dos alunos ( Gráfico 4.3, Gráfico 4.4, e, no Anexo II, a Tabela II.2 ).
A análise das respostas dos estudantes quanto ao grau de dificuldade do Componente de Conhecimento Específico da prova, agregado por Grande Região, mostra que a diferença entre a maior e a menor proporção de alunos que a avaliaram como difícil ou muito difícil é estatisticamente significativa : 55,6 \% na região Norte e 13,0 \% na Centro-Oeste.
O percentual de alunos que classificaram o grau de dificuldade como médio, no Componente de Conhecimento Específico, variou de 44,4 \% a 87,0 \%, valores referentes às regiões Norte e Centro-Oeste, respectivamente.
Nenhum estudante nas regiões Norte e Centro-Oeste avaliou a prova como fácil ou muito fácil.
Considerando a avaliação da dificuldade das questões do Componente de Conhecimento Específico da prova, de acordo com o desempenho dos estudantes, não se observa diferença estatisticamente significativa, a não ser entre os dois quartos extremos.
Em todos os quartos, a proporção dos que classificaram a parte específica como difícil ou muito difícil variou de 11,8 \% ( 4º quarto ) a 26,6 \% ( 1º quarto ), sendo tanto menor quanto maior fosse o quarto de desempenho.
A alternativa modal para a Questão 2 foi o médio, com 62,4 \% no dois quartos inferiores, e 71,8 \% do estudantes do quarto superior optando por esta resposta. 4.2 EXTENSÃO DA PROVA EM RELAÇÃO AO TEMPO TOTAL Indagados quanto à extensão da prova, em relação ao tempo total oferecido para a sua resolução ( Questão 3 ), os estudantes apontaram, com maior incidência, a alternativa que considerava a extensão adequada, para todas as agregações consideradas ( Gráfico 4.5, Gráfico 4.6, e, no Anexo II, a Tabela II.3 ).
Em termos nacionais, o percentual de alunos que responderam ser a extensão da prova adequada foi de 57,8 \%.
Já 34,3 \% dos inscritos presentes consideraram que o exame foi longo ou muito longo e menos do que 8 \% o avaliaram como curto ou muito curto.
Entre as Grandes Regiões a proporção daqueles que avaliaram a prova como longa ou muito longa em relação ao tempo total destinado à sua resolução variou pouco : de 25,9 \% na região Norte até 39,1 \% na região Centro-Oeste.
As diferenças entre as regiões não são estaticamente significativas.
Observando o desempenho dos alunos, nota -se que a proporção de estudantes que considerou a extensão da prova adequada é crescente com o desempenho dos estudantes.
Variou de 47,7 \% no 1º quarto a 65,5 no 4º quarto.
No Gráfico 4.6, pode -se constatar que as proporções dos estudantes que consideraram a prova longa ou muito longa, em relação ao tempo total destinado à sua resolução, decrescem com os quartos de desempenho, e as diferenças entre os quartos não são estatisticamente significativas. 4.3 COMPREENSÃO DOS ENUNCIADOS DAS QUESTÕES 4.3.1 Componente de Formação Geral Com relação aos enunciados das questões do Componente de Formação Geral ( Questão 4 ), as opiniões foram positivas, já que 72,9 \% dos alunos avaliados consideraram os enunciados de todas ou da maioria das questões claros e objetivos ( Gráfico 4.7, Gráfico 4.8, e, no Anexo II, a Tabela II.4 ).
Na análise regional, a percentagem de estudantes que avaliaram que todos ou a maioria dos enunciados das questões do Componente de Formação Geral estavam claros e objetivos variou de 64,3 \% na região Sul a 82,6 \% na região Centro-Oeste, não ocorrendo diferenças estatisticamente significativas entre as Grandes Regiões.
A análise das percepções dos estudantes sobre a clareza e objetividade dos enunciados permite afirmar que todos ou a maioria dos enunciados de questões relativas ao Componente de Formação Geral foram considerados claros e objetivos para a maior parte dos respondentes ( maior do que 64 \% em todas as regiões e maior do que 65 \% em todos os quartos de desempenho ).
Segundo o desempenho, observa -se que a proporção dos que os enunciados de todas ou da maioria das questões claros e objetivos cresce conforme o desempenho aumenta, com diferenças estatisticamente não significativas.
No quarto superior, a clareza e objetividade de todos ou da maioria dos enunciados das questões foi percebida por 79,1 \% dos alunos e no quarto de desempenho inferior tal avaliação foi emitida por 65,1 \% deles. 4.3.2 Componente de Conhecimento Específico Com relação aos enunciados das questões do Componente de Conhecimento Específico da prova, para 78,1 \% dos estudantes avaliados da Área de Tecnologia em Construção de Edifícios a clareza e a objetividade ( Questão 5 ) estavam presentes em todas ou na maioria das questões ( Gráfico 4.9, Gráfico 4.10, e no Anexo II, a Tabela II.5 ).
A maioria dos estudantes de todas as Grandes Regiões brasileiras considerou claros e objetivos todos ou a maioria dos enunciados das questões do Componente de Conhecimento Específico da prova, percentual sempre maior do que 74 \%.
As diferenças entre os percentuais das Grandes Regiões não são estatisticamente significativa.
A proporção de estudantes que consideraram os enunciados das questões claros e objetivos assume o menor valor no quarto inferior de desempenho ( 71,6 \% ), e o maior valor no quarto superior ( 81,8 \% ).
As diferenças entre as proporções dos quartos de desempenho não são estatisticamente significativas. 4.4 SUFICIÊNCIA DAS INFORMAÇÕES/INSTRUÇÕES FORNECIDAS Ao avaliarem as informações/instruções fornecidas para a resolução das questões ( Questão 6 ), 85,3 \% dos respondentes da Área de Tecnologia em Construção de Edifícios de todo o Brasil afirmaram que estas eram até excessivas ou suficientes em todas ou na maioria das questões ( Gráfico 4.11, Gráfico 4.12, e, no Anexo II, a Tabela II.6 ).
Quanto à distribuição de respondentes pelas Grandes Regiões observa -se que a proporção de estudantes que consideraram as informações/instruções fornecidas até excessivas ou suficientes em todas ou na maioria das questões foi sempre superior a 80 \%, chegando a 92,9 \% na região Sul.
As diferenças entre as regiões não são estatisticamente significativas.
Levando -se em conta o desempenho dos participantes, entre os que responderam que as informações/instruções fornecidas para a resolução das questões eram até excessivas ou suficientes em todas ou na maioria das questões, somente os quartos extremos apresentam diferença estatisticamente significativa, como mostra o Gráfico 4.12.
O percentual de participantes que apresentaram tal avaliação foi crescente com o aumento do desempenho : 79,8 \% dos respondentes no quarto inferior, e 92,7 \% no quanto superior. 4.5 DIFICULDADE ENCONTRADA AO RESPONDER À PROVA Perguntados sobre as dificuldades com as quais se depararam ao responder à prova ( Questão 7 ), 10,0 \% dos estudantes apontaram o desconhecimento do conteúdo.
Para 39,4 \%, a forma diferente de abordagem do conteúdo foi indicada como dificuldade.
Já a falta de motivação para fazer a prova foi a dificuldade apontada por 21,2 \% dos respondentes.
Considerando -se todo o Brasil, 21,4 \% dos respondentes afirmaram que não tiveram qualquer tipo de dificuldade para responder à prova ( Tabela II.7 no Anexo II ).
Os Gráficos 4.13 e 4.14 apresentam os percentuais de estudantes que apontaram o desconhecimento do conteúdo como dificuldade percebida ao responder à prova.
Na análise por Grandes Regiões, o percentual de inscritos e presentes que apontaram o desconhecimento do conteúdo como dificuldade ao responder à prova não superou 16 \%.
Os percentuais variaram de 4,3 \% na região Centro-Oeste a 15,4 \% na Sul.
O Gráfico 4.13 revela que não há diferenças significativas estatisticamente entre as Grandes Regiões..
A forma diferente de abordagem do conteúdo foi a escolha modal dos estudantes, excetuando -se os estudantes da região Centro-Oeste que tiveram como escolha modal a falta de motivação para fazer a prova.
O percentual de alunos que citaram a falta de motivação como dificuldade variou de 0,0 \% ( região Sul ) a 39,1 \% ( região Centro-Oeste ).
Os que declararam não ter qualquer dificuldade para responder à prova variaram de 11,1 \% na região Norte a 30,4 \% na Centro-Oeste.
Com relação aos quartos de desempenho, o desconhecimento do conteúdo, foi a opção escolhida por 15,5 \% dos estudantes do quarto superior e por 9,3 \% do quarto inferior.
A alternativa modal para os alunos, quando agregados pelos quartos de desempenho, foi a forma diferente de abordagem do conteúdo : 36,4 \% no quarto inferior e 33,6 \% do quarto superior assim o responderam. 4.6 CONTEÚDOS DAS QUESTÕES OBJETIVAS DA PROVA Ao analisarem os conteúdos das questões objetivas da prova ( Questão 8 ), um percentual muito pequeno dos estudantes avaliados, apenas 3,9 \%, afirmou que não estudou ainda a maioria desses conteúdos ( Gráficos 4.15, Gráfico 4.16, e a Tabelas II.8 no Anexo II ).
A maioria ( 77,7 \% ) afirmou ter estudado e aprendido muitos ou todos os conteúdos avaliados.
Na análise por Grande Região, a proporção de respondentes que escolheram a opção não estudou ainda a maioria desses conteúdos foi pequena, sendo que nenhum estudante na região Centro-Oeste fez essa escolha.
Na região Sudeste ( 2,1 \% ), apesar de pequena, a proporção foi menor do que a média nacional ( 3,9 \% ).
Não se observa diferença estatisticamente significativa entre as regiões.
Em todas as regiões, a maioria dos presentes afirmou ter estudado e aprendido muitos ou todos os conteúdos, com proporções variando entre 66,7 \% na região Norte a 91,3 \% na Centro-Oeste.
Considerando -se separadamente as opiniões de estudantes dos quatro quartos de desempenho, observa -se que, no quarto inferior, 9,4 \% ofereceram como resposta que não estudou ainda a maioria desses conteúdos, sendo menos de 2,0 \% os do quarto superior com a mesma resposta.
As diferenças entre os percentuais dos alunos que optaram por este motivo de dificuldade não são estatisticamente significativa.
Tendo em conta o quarto superior, 90,0 \% dos alunos afirmaram ter estudado e aprendido muitos ou todos os conteúdos. 4.7 TEMPO GASTO PARA CONCLUIR A PROVA Ao responderem sobre o tempo de conclusão da prova ( Questão 9 ), mais de dois terços dos estudantes ( 74,7 \% ) afirmaram ter gasto entre duas e quatro horas ( Gráfico 4.17, Gráfico 4.18 e, no Anexo II, a Tabela II.9 ).
Considerando as cinco Grandes Regiões brasileiras, os que utilizaram entre duas e quatro horas para finalizar a prova ficaram abaixo do percentual nacional na região Norte ( 70,4 \% ) e na Nordeste ( 73,7 \% ).
Entre as regiões não existem diferenças estatisticamente significativas, como mostra o Gráfico 4.17.
Uma vez sendo analisadas as alternativas escolhidas pelos estudantes que se situam nos diferentes quartos de desempenho, nota -se que a maior proporção de participantes que declarou ter gasto entre duas e quatro horas para concluir a prova encontra -se no quarto superior ( 84,5 \% ), enquanto a menor proporção é observada no quarto inferior ( 55,0 \% ).
As diferenças entre o 2º, 3º e 4º quartos não são estatisticamente significativas.

% matched lemmas: 
\end{textsample}
